% ===== 原文 pp.17–18（PDF pp.18–19）=====
% tr-02 片段 a：§4 标题与引言、4.1 数据、4.2 拉氏量（开头）
% 注：本片段自原 PDF p.18 中部开始（该页上半是 §3.3 结尾，属tr-01）；
%     印刷 p.17 的页标记 % [原文 p.17] 已由 tr-01 片段给出。
\section{规范线性 sigma 模型}

本节介绍一类二维$(2,2)$超对称规范理论，称为规范线性sigma模型（GLSM），它由[11]引入。

对紧李群$G$，我们分别用$T$和$Z_G$表示一个极大环面与中心。我们用$\mathrm{W}$表示$G$的Weyl群。我们把$\fg\supset\ft\supset\fz$记作$G\supset T\supset Z_G$的李代数，并把它们视为“纯虚”的。“实”的则是复化李代数$\fg_\C\supset\ft_\C\supset\fz_\C$中的$i\fg\supset i\ft\supset i\fz$。于是，对取值于$\fg_\C$的量$X$，我们写成$X=\Re(X)+i\Im(X)$，其中$\Re(X)$与$\Im(X)$都取值于$i\fg$，并记$\overline{X}=\Re(X)-i\Im(X)$。$T$的权格记为$\mathrm{P}\subset i\ft^*$。

% [Tutor Note: p.17 的“纯虚”约定不是笔误：把$\fg,\ft,\fz$取为纯虚后，$i\fg$才是实李代数，(4.8) 的双线性型$(X,Y)\in i\fg\times i\fg\mapsto\frac{1}{e^2}XY\in\R$才能取实值（否则$XY$为负定，无法作实内积）。由 (4.8) 的像落在$\R$即可反推$X,Y\in i\fg$；这也是 (4.10)、(4.11) 中$D$、$v_{01}$称为“实值”而$\sigma,\lambda$取值于$\fg_\C$的原因。]

\subsection{数据}

一个二维$(2,2)$规范理论由如下选择确定：

\begin{itemize}
  \item \emph{规范群}$G$：一个紧李群；
  \item \emph{物质表示}$V$：$G$的一个有限维复表示；
  \item \emph{超势}$W(\phi)$：$\phi\in V$的一个$G$不变多项式函数；以及
  \item \emph{扭超势}$\widetilde{W}(\sigma)$：$\sigma\in\fg_\C$的一个$G$不变多项式函数。
\end{itemize}

作为数据中较次要的部分，我们还选取$i\fg$上的一个$G$不变内积与$V$上的一个$G$不变Hermite内积，二者均为正定。

% [原文 p.18]

当存在与$G$-作用交换的线性映射$R:V\to V$，使得
\begin{equation}
W(\lambda^R\phi)=\lambda^2W(\phi)
\tag{4.1}
\end{equation}
时，便存在一个矢量$U(1)$ R-对称性。

当$\ee^{\pi iR}:V\to V$与某个规范群元素（记作$J\in G$）的作用相同时，荷整性（charge integrality）成立。当$\widetilde{W}(\sigma)$为线性时，经典层面存在一个带有荷整性的轴$U(1)$ R-对称性；而在Calabi–Yau条件$G\subset SL(V)$之下，它仍然是量子系统的对称性。在本文中，除Calabi–Yau条件之外，我们假定上述一切都成立。我们把线性扭超势写成
\begin{equation}
\widetilde{W}(\sigma)=-t(\sigma),
\tag{4.2}
\end{equation}
其中
\begin{equation}
t=\zeta-i\theta\in\fg^{*G}_\C,
\tag{4.3}
\end{equation}
是一个伴随不变线性形式；这里$\zeta$与$\theta$（二者都在$i\fg^{*G}$中）分别称为Fayet-Illiopoulos（FI）参数与theta参数。注意，借助自然同构$\fg^{*G}\cong\ft^{*W}\cong\fz^*$，$\zeta$与$\theta$也可视为$i\ft^{*W}$或$i\fz^*$中的元素。确切地说，theta参数要服从一个离散等同
\begin{equation}
\theta\equiv\theta+2\pi n,
\tag{4.4}
\end{equation}
其中$n$是某个特征$G\to U(1)$在微分映射$\Hom(G,U(1))\to\Hom(\fg,\fu(1))=i\fg^*$下的像。因此，theta参数（或theta角）的空间是紧环面$i\fg^{*G}/2\pi\Lambda_G$，其中$\Lambda_G:=\Image(\Hom(G,U(1))\to i\fg^{*G})$。当$G$连通时，$\Lambda_G$同构于$\Hom(G,U(1))$，并且等于由$i\ft^{*W}\cong i\fg^{*G}$嵌入$i\fg^{*G}$的、$T$的Weyl不变权所构成的格$\mathrm{P}^W$。

% [Tutor Note: p.18 式 (4.3) 的符号$t=\zeta-i\theta$必须与 (4.2)、(4.13) 联读：由 (4.2) $\widetilde{W}=-t$及 (4.9) 中$\Re Q_+\overline{Q}_-(-t(\sigma))$可知，$\zeta$部分给出 (4.13) 的$-\zeta(D)$，$\theta$部分给出$+\theta(v_{01})$；若把 (4.3) 抄成$\zeta+i\theta$，(4.13) 的$\theta$项符号就会反过来。另：原文把 Fayet–Iliopoulos 拼作“Fayet-Illiopoulos”，此处照录。]

\subsection{拉氏量}

该理论由两组场构成，分别称为物质多重态与规范多重态。规范多重态由一个$G$联络$v_\mu$，以及一个标量$\sigma$、一个Dirac费米子$\lambda$和一个标量$D$组成，它们分别取值于$\fg_\C$、$\fg_\C$与$i\fg$。物质多重态由一个标量$\phi$、一个Dirac费米子$\psi$和一个标量$F$组成，三者都取值于$V$。规范联络$v_\mu$是“实值”的，因此曲率为$v_{\mu\nu}=\partial_\mu v_\nu-\partial_\nu v_\mu+i[v_\mu,v_\nu]$，而规范协变导数为$D_\mu\phi=\partial_\mu\phi+iv_\mu\phi$。

我们在Minkowski空间上描述超对称与拉氏量，其时间与空间坐标分别为$x^0$与$x^1$，度规为$\dd s^2=-(\dd x^0)^2+(\dd x^1)^2$。我们经常使用光锥坐标$x^\pm=x^0\pm x^1$。
